\chapter{Limites de l'étude}

Cette section est volontairement distincte de la synthèse
(chapitre~11) : elle ne redonne pas de résultats, elle cadre ce que ces
résultats permettent — et ne permettent pas — d'affirmer.

\section{Le dataset ne permet pas d'affirmer de causalité}

Aucune conclusion de ce rapport ne doit être lue comme \og l'IA cause
X sur l'emploi \fg{}. Le dataset décrit un ensemble d'offres à un
instant donné (ou sur une période donnée) ; il ne contient ni série
longitudinale par individu, ni groupe de contrôle, ni variable
d'intervention. Les corrélations calculées au chapitre~9 §9.3.3
(salaire/expérience, salaire/demande) sont des associations
statistiques descriptives, pas des relations causales — et restent de
toute façon faibles dans l'absolu (aucune ne dépasse 0,3).

\section{Représentativité incertaine du dataset}

\begin{itemize}[leftmargin=1.5em]
  \item \textbf{Volume limité} : 1500 offres, 25 métiers, 14 pays — un
    échantillon, pas un recensement exhaustif du marché mondial de
    l'emploi IA.
  \item \textbf{Absence totale d'anomalie} : 0 valeur manquante sur les
    25 colonnes, 0 doublon, 0 salaire invalide (chapitre~2 §2.3,
    chapitre~6) — une propreté inhabituelle pour une extraction de
    données réelle, qui suggère un dataset au moins partiellement
    généré/synthétisé plutôt que collecté brut. Cela ne remet pas en
    cause les calculs effectués, mais invite à ne pas sur-interpréter
    la précision apparente des chiffres comme une exactitude du marché
    réel.
  \item \textbf{Répartition d'expérience quasi uniforme} (Entry 385,
    Mid 370, Senior 364, Lead 381 — chapitre~2 §2.3) : un marché réel
    présente rarement une distribution aussi équilibrée entre niveaux
    d'expérience, ce qui renforce l'hypothèse d'un jeu de données
    construit pour l'exercice pédagogique plutôt qu'une extraction
    représentative d'un marché réel.
  \item \textbf{Provenance Kaggle non vérifiable} : la méthode de
    collecte du dataset \emph{AI Jobs Market 2025-2026 | Salaries}
    (scraping, enquête, génération synthétique) n'est pas documentée
    par sa source — impossible de garantir qu'il reflète fidèlement les
    pratiques salariales réelles du secteur.
\end{itemize}

\section{Limites méthodologiques spécifiques}

\begin{itemize}[leftmargin=1.5em]
  \item \textbf{Pic d'activité 2026 (constat~8, chapitre~11)} : la
    comparaison janvier-mars 2026 (271 à 306 offres/mois) contre
    l'ensemble de 2025 (41 à 68/mois) oppose 3 mois à 12 mois — une
    rupture de volume aussi nette peut refléter une réelle accélération
    du marché, ou un artefact de construction du dataset (échantillonnage
    non uniforme dans le temps, sur-représentation délibérée de la
    période récente). Le dataset ne permet pas de trancher entre ces
    deux hypothèses.
  \item \textbf{Corrélations de Pearson uniquement} (chapitre~9
    §9.3.3) : \texttt{CORR()} mesure une relation \emph{linéaire}
    entre variables continues ; elle peut sous-estimer une relation
    réelle non linéaire (effet de seuil par tranche catégorielle,
    observé par ailleurs sur \texttt{experience\_level} au KPI~7) sans
    que cela contredise le résultat obtenu, comme discuté au chapitre~9.
  \item \textbf{\texttt{ai\_salary\_premium\_pct} sans référentiel
    documenté} (KPI~9, chapitre~9) : le dataset ne précise pas contre
    quel salaire de référence ce premium est calculé — les
    comparaisons entre métiers ou catégories restent valides
    relativement les unes aux autres, mais leur valeur absolue ne peut
    pas être vérifiée indépendamment.
  \item \textbf{Un seul dataset, une seule source} : aucune comparaison
    à une source externe (LinkedIn, Glassdoor, INSEE) n'a été réalisée
    dans ce rapport (piste identifiée comme bonus, voir issue~\#116,
    hors périmètre du rendu principal) — les constats du chapitre~11
    valent pour ce dataset, pas comme vérité générale du marché de
    l'emploi IA.
\end{itemize}
